\documentclass[11pt]{article}
\usepackage{amsmath, amssymb, amsfonts, amsthm}
\parindent 0px

\begin{document}

\setcounter{section}{2}
\section{Transformations for 3D movement}

We can break down $\vec{P_{3D}}$ into 2 states: $\vec{P_{3DI}}$ and $\vec{P_{3DF}}$.
$\vec{P_{3DI}}$ represents the location of a point before a set of transformations,
$\{ T_{0}, T_{1}, \ldots, T_{n} \}$, that describe its movement. $\vec{P_{3DF}}$ 
represents the location of a point after the set of transformations.\\

Let $T_{n}:\mathbb{R}^3 \rightarrow \mathbb{R}^3$.\\
Let $T_{n}$ be of the form 
$\left\{ T_{n}(step) \left| \vec{P_{3DF}} = 
 \mathbf{A}\vec{P_{3DI}} + \vec{b}\right. \right\}$ \\

Keypresses are polled for 1 or more of the following transformations\\
$\left\{ T_{rx}(\Delta\theta), T_{ry}(\Delta\theta), T_{rz}(\Delta\theta),
T_z(\Delta z) \right\}$, which apply to all points during each iteration of the
render loop. The table below describes 
the function and expression of each transformation.\\
\\

\noindent\hspace*{-2.7cm}
\begin{tabular}{ |c|c|c|c| }
	\hline
	Our POV & Point Movement & $T_n$ & Expression \\
	\hline\hline

	\rule{0pt}{3.0ex} Move forward & Lower Z-component value & $T_z(-\Delta z)$ &
	$\vec{P_{3DF}} = \mathbf{I_3}\vec{P_{3DI}} - \Delta z \hat{z}$ \\
	\hline

	\rule{0pt}{3.0ex} Move backward & Raise Z-component value & $T_z(+\Delta z)$ &
	$\vec{P_{3DF}} = \mathbf{I_3}\vec{P_{3DI}} + \Delta z \hat{z}$ \\
	\hline

	\rule{0pt}{3.0ex} Pitch up & Rotate around x-axis counter-clockwise &
	$T_{rx}(+\Delta\theta)$ &
	$\vec{P_{3DF}} = \mathbf{R_x}(+\Delta\theta)\vec{P_{3DI}} + \vec{P_{proj}}$ \\
	\hline

	\rule{0pt}{3.0ex} Pitch down & Rotate around x-axis clockwise &
	$T_{rx}(-\Delta\theta)$ &
	$\vec{P_{3DF}} = \mathbf{R_x}(-\Delta\theta)\vec{P_{3DI}} + \vec{P_{proj}}$ \\
	\hline

	\rule{0pt}{3.0ex} Yaw left & Rotate around y-axis counter-clockwise &
	$T_{ry}(+\Delta\theta)$ &
	$\vec{P_{3DF}} = \mathbf{R_y}(+\Delta\theta)\vec{P_{3DI}} + \vec{P_{proj}}$ \\
	\hline

	\rule{0pt}{3.0ex} Yaw right & Rotate around y-axis clockwise &
	$T_{ry}(-\Delta\theta)$ &
	$\vec{P_{3DF}} = \mathbf{R_y}(-\Delta\theta)\vec{P_{3DI}} + \vec{P_{proj}}$ \\
	\hline

	\rule{0pt}{3.0ex} Roll left & Rotate around z-axis clockwise &
	$T_{rz}(-\Delta\theta)$ &
	$\vec{P_{3DF}} = \mathbf{R_z}(-\Delta\theta)\vec{P_{3DI}} + \vec{P_{proj}}$ \\
	\hline

	\rule{0pt}{3.0ex} Roll right & Rotate around z-axis counter-clockwise &
	$T_{rz}(+\Delta\theta)$ &
	$\vec{P_{3DF}} = \mathbf{R_z}(+\Delta\theta)\vec{P_{3DI}} + \vec{P_{proj}}$ \\
	\hline

\end{tabular}
\\
\\

\textbf{Expansions:}
$$\vec{P_{proj}} = \begin{bmatrix} 0\\ 0\\ 0 \end{bmatrix}$$
$$\hat{z} = \begin{bmatrix} 0\\ 0\\ 1 \end{bmatrix}$$

$$\mathbf{I_3} = \begin{bmatrix}
				1 & 0 & 0 \\
				0 & 1 & 0 \\
				0 & 0 & 1
		 \end{bmatrix}$$

$$\mathbf{R_x}(\theta) = \begin{bmatrix}
				        1 &        0      &       0      \\
				        0 &  \cos(\theta) & \sin(\theta) \\
				        0 & -\sin(\theta) & \cos(\theta)
			\end{bmatrix}$$

$$\mathbf{R_y}(\theta) = \begin{bmatrix}
				        \cos(\theta) & 0 & -\sin(\theta) \\
				              0      & 1 &        0	 \\
				        \sin(\theta) & 0 &  \cos(\theta)
			\end{bmatrix}$$

$$\mathbf{R_z}(\theta) = \begin{bmatrix}
				         \cos(\theta) & \sin(\theta) & 0 \\
				        -\sin(\theta) & \cos(\theta) & 0 \\
				              0       &       0      & 1
			\end{bmatrix}$$


\end{document}

